% GENERATED FILE. Edit canonical lessons, then run npm run generate:books.
% canonical-source-hash: fnv1a64:f3af803669289847
% canonical-lessons: GU-C300-angaro, GU-C300-dhumado, GU-C300-jvala, GU-C300-salgavvu, GU-C300-dajhvu, GU-R300-review-293-300

\chapter{Smoke from the Kitchen}
\label{ch:gu-smoke-from-the-kitchen}
\hlchaptermodality{300}

\begin{chapteropening}
\textbf{By the end of this chapter:} \emph{I can say how a fire started and who got burnt.}

Complete the last lesson of chapter 300: I can say how a fire started and who got burnt.
\end{chapteropening}

\section[\texorpdfstring{\gu{અંગારો} (angāro)}{અંગારો (angāro)}]{\gu{અંગારો} (angāro) — an ember}

\label{lesson:GU-C300-angaro}



\begin{quote}
Before the new one: say the Gujarati for to hold back, then the Gujarati for safe.
\end{quote}

\subsection*{You'll want to know: \gu{અંગારો}}
\textbf{\gu{અંગારો}} — \emph{angāro} — \textquotedblleft{}an ember\textquotedblright{}.

\begin{quote}
\gu{ચૂલામાં} \gu{હજી} \gu{અંગારા} \gu{સળગતા} \gu{હતા}.
\end{quote}

\begin{quote}
\emph{cūlāmā̃ hajī angārā saḷagtā hatā.}
\end{quote}

\begin{quote}
Embers were still glowing in the stove.
\end{quote}

A masculine noun.

\begin{cousinweb}
\textbf{\gu{અંગારો}} (\emph{angāro}) is inherited from Sanskrit \emph{aṅgāra}, \textquotedblleft{}charcoal\textquotedblright{}, and CDIAL lists Gujarati \emph{ãgāro}, \textquotedblleft{}glowing charcoal\textquotedblright{} (CDIAL 125). Embers stay in a stove after the flames have gone: \textbf{\gu{અંગારા}}.
\end{cousinweb}

\begin{grammarlens}[title={What you've built}]
One more word for telling what happened.
\end{grammarlens}

\subsection*{Your turn}
\begin{itemize}
\raggedright
  \item \emph{Say it:} \emph{angāro}
  \item \emph{Say it:} \emph{angāro}, once more
  \item \emph{Say it:} \emph{cūlāmā̃ hajī angārā saḷagtā hatā.} — Embers were still glowing in the stove.
  \item \emph{Recall:} say the Gujarati for to hold back, then the Gujarati for safe, then say \emph{angāro} again
\end{itemize}

\begin{tcolorbox}[breakable,colback=teal!4,colframe=teal!35!black,title={Before you move on}]
What does \gu{અંગારો} mean? (\textquotedblleft{}An ember\textquotedblright{}.) Say it once more.
\end{tcolorbox}
\section[\texorpdfstring{\gu{ધુમાડો} (dhumāḍo)}{ધુમાડો (dhumāḍo)}]{\gu{ધુમાડો} (dhumāḍo) — smoke}

\label{lesson:GU-C300-dhumado}



\begin{quote}
Before the new one: say the Gujarati for safe, then the Gujarati for an ember.
\end{quote}

\subsection*{You'll want to know: \gu{ધુમાડો}}
\textbf{\gu{ધુમાડો}} — \emph{dhumāḍo} — \textquotedblleft{}smoke\textquotedblright{}.

\begin{quote}
\gu{રસોડામાંથી} \gu{કાળો} \gu{ધુમાડો} \gu{નીકળતો} \gu{હતો}.
\end{quote}

\begin{quote}
\emph{rasoḍāmā̃thī kāḷo dhumāḍo nīkaḷto hato.}
\end{quote}

\begin{quote}
Black smoke was coming out of the kitchen.
\end{quote}

A masculine noun.

\begin{cousinweb}
\textbf{\gu{ધુમાડો}} (\emph{dhumāḍo}) is inherited from Sanskrit \emph{dhūmá}, \textquotedblleft{}smoke\textquotedblright{}, and CDIAL lists the Gujarati word (CDIAL 6849). It is masculine: \textbf{\gu{કાળો} \gu{ધુમાડો}}, black smoke.
\end{cousinweb}

\begin{grammarlens}[title={What you've built}]
One more word for telling what happened.
\end{grammarlens}

\subsection*{Your turn}
\begin{itemize}
\raggedright
  \item \emph{Say it:} \emph{dhumāḍo}
  \item \emph{Say it:} \emph{dhumāḍo}, once more
  \item \emph{Say it:} \emph{rasoḍāmā̃thī kāḷo dhumāḍo nīkaḷto hato.} — Black smoke was coming out of the kitchen.
  \item \emph{Recall:} say the Gujarati for safe, then the Gujarati for an ember, then say \emph{dhumāḍo} again
\end{itemize}

\begin{tcolorbox}[breakable,colback=teal!4,colframe=teal!35!black,title={Before you move on}]
What does \gu{ધુમાડો} mean? (\textquotedblleft{}Smoke\textquotedblright{}.) Say it once more.
\end{tcolorbox}
\section[\texorpdfstring{\gu{જ્વાળા} (jvāḷā)}{જ્વાળા (jvāḷā)}]{\gu{જ્વાળા} (jvāḷā) — a flame}

\label{lesson:GU-C300-jvala}



\begin{quote}
Before the new one: say the Gujarati for an ember, then the Gujarati for smoke.
\end{quote}

\subsection*{You'll want to know: \gu{જ્વાળા}}
\textbf{\gu{જ્વાળા}} — \emph{jvāḷā} — \textquotedblleft{}a flame\textquotedblright{}.

\begin{quote}
\gu{આગની} \gu{જ્વાળાઓ} \gu{છત} \gu{સુધી} \gu{પહોંચી}.
\end{quote}

\begin{quote}
\emph{āgnī jvāḷāo chat sudhī pahõcī.}
\end{quote}

\begin{quote}
The flames reached the ceiling.
\end{quote}

A feminine noun.

\begin{cousinweb}
\textbf{\gu{જ્વાળા}} (\emph{jvāḷā}) is a tatsama, Sanskrit \emph{jvālā}, \textquotedblleft{}a flame, a blaze\textquotedblright{} (Monier-Williams). It is feminine, and fire comes in the plural: \textbf{\gu{જ્વાળાઓ}}, flames.
\end{cousinweb}

\begin{grammarlens}[title={What you've built}]
One more word for telling what happened.
\end{grammarlens}

\subsection*{Your turn}
\begin{itemize}
\raggedright
  \item \emph{Say it:} \emph{jvāḷā}
  \item \emph{Say it:} \emph{jvāḷā}, once more
  \item \emph{Say it:} \emph{āgnī jvāḷāo chat sudhī pahõcī.} — The flames reached the ceiling.
  \item \emph{Recall:} say the Gujarati for an ember, then the Gujarati for smoke, then say \emph{jvāḷā} again
\end{itemize}

\begin{tcolorbox}[breakable,colback=teal!4,colframe=teal!35!black,title={Before you move on}]
What does \gu{જ્વાળા} mean? (\textquotedblleft{}A flame\textquotedblright{}.) Say it once more.
\end{tcolorbox}
\section[\texorpdfstring{\gu{સળગાવવું} (saḷgāvvũ)}{સળગાવવું (saḷgāvvũ)}]{\gu{સળગાવવું} (saḷgāvvũ) — to set alight; to light}

\label{lesson:GU-C300-salgavvu}



\begin{quote}
Before the new one: say the Gujarati for smoke, then the Gujarati for a flame.
\end{quote}

\subsection*{You'll want to know: \gu{સળગાવવું}}
\textbf{\gu{સળગાવવું}} — \emph{saḷgāvvũ} — \textquotedblleft{}to set alight; to light\textquotedblright{}.

\begin{quote}
\gu{કોઈએ} \gu{કચરો} \gu{સળગાવ્યો} \gu{અને} \gu{આગ} \gu{ફેલાઈ} \gu{ગઈ}.
\end{quote}

\begin{quote}
\emph{koīe kacro saḷgāvyo ane āg phelāī gaī.}
\end{quote}

\begin{quote}
Someone set fire to the rubbish and the fire spread.
\end{quote}

Past, with \gu{તેણે} (he or she …): \gu{તેણે} \gu{ચૂલો} \gu{સળગાવ્યો}, \emph{teṇe cūlo saḷgāvyo}.

\begin{cousinweb}
\textbf{\gu{સળગાવવું}} (\emph{saḷgāvvũ}) is the causative of \textbf{\gu{સળગવું}}, to catch fire, which goes back to Sanskrit \emph{saṁlagna}, \textquotedblleft{}closely attached\textquotedblright{}; CDIAL lists both verbs (CDIAL 12999). You light a stove with it too: \textbf{\gu{ચૂલો} \gu{સળગાવ્યો}}.
\end{cousinweb}

\begin{grammarlens}[title={What you've built}]
One more word for telling what happened.
\end{grammarlens}

\subsection*{Your turn}
\begin{itemize}
\raggedright
  \item \emph{Say it:} \emph{saḷgāvvũ}
  \item \emph{Say it:} \emph{saḷgāvvũ}, once more
  \item \emph{Say it:} \emph{koīe kacro saḷgāvyo ane āg phelāī gaī.} — Someone set fire to the rubbish and the fire spread.
  \item \emph{Recall:} say the Gujarati for smoke, then the Gujarati for a flame, then say \emph{saḷgāvvũ} again
\end{itemize}

\begin{tcolorbox}[breakable,colback=teal!4,colframe=teal!35!black,title={Before you move on}]
What does \gu{સળગાવવું} mean? (\textquotedblleft{}To set alight; to light\textquotedblright{}.) Say it once more.
\end{tcolorbox}
\section[\texorpdfstring{\gu{દાઝવું} (dājhvũ)}{દાઝવું (dājhvũ)}]{\gu{દાઝવું} (dājhvũ) — to be scalded; to get a burn}

\label{lesson:GU-C300-dajhvu}



\begin{quote}
Before the new one: say the Gujarati for a flame, then the Gujarati for to set alight.
\end{quote}

\subsection*{You'll want to know: \gu{દાઝવું}}
\textbf{\gu{દાઝવું}} — \emph{dājhvũ} — \textquotedblleft{}to be scalded; to get a burn\textquotedblright{}.

\begin{quote}
\gu{ગરમ} \gu{તેલથી} \gu{મીરાનો} \gu{હાથ} \gu{દાઝી} \gu{ગયો}.
\end{quote}

\begin{quote}
\emph{garam telthī Mīrāno hāth dājhī gayo.}
\end{quote}

\begin{quote}
Mira burned her hand on hot oil.
\end{quote}

Past (he …): \gu{તે} \gu{છોકરો} \gu{દાઝી} \gu{ગયો}, \emph{te chokro dājhī gayo}.

\begin{cousinweb}
\textbf{\gu{દાઝવું}} (\emph{dājhvũ}) is inherited from Sanskrit \emph{dahyátē}, \textquotedblleft{}is burnt\textquotedblright{}, and CDIAL lists the Gujarati verb, \textquotedblleft{}to be burnt, be scalded\textquotedblright{} (CDIAL 6248). It is skin that \textbf{\gu{દાઝે} \gu{છે}}; wood and paper \textbf{\gu{બળે} \gu{છે}}.
\end{cousinweb}

\begin{grammarlens}[title={What you've built}]
One more word for telling what happened.
\end{grammarlens}

\subsection*{Your turn}
\begin{itemize}
\raggedright
  \item \emph{Say it:} \emph{dājhvũ}
  \item \emph{Say it:} \emph{dājhvũ}, once more
  \item \emph{Say it:} \emph{garam telthī Mīrāno hāth dājhī gayo.} — Mira burned her hand on hot oil.
  \item \emph{Recall:} say the Gujarati for a flame, then the Gujarati for to set alight, then say \emph{dājhvũ} again
\end{itemize}

\begin{tcolorbox}[breakable,colback=teal!4,colframe=teal!35!black,title={Before you move on}]
What does \gu{દાઝવું} mean? (\textquotedblleft{}To be scalded; to get a burn\textquotedblright{}.) Say it once more.
\end{tcolorbox}
\section[(dialogue)]{Looking back: from An Accident near the Bridge to Smoke from the Kitchen}

\label{lesson:GU-R300-review-293-300}



\begin{quote}
Say the Gujarati for an accident and a collision.
\end{quote}

\subsection*{Your turn}
Say these aloud:

\begin{itemize}
\raggedright
  \item an accident, a collision, injured, serious, to panic
  \item to slip, a graze, to tumble, a slope, consciousness
  \item a bend, a brake, careless, a fine, to crush
  \item an ambulance, to pick up, a bandage, a stitch, a condition
  \item to warn, careful, upside down, a gathered crowd, a push
  \item death, mourning, a shock, consolation, to bear
  \item a scream, a shout, to bring to a stop, to hold back, safe
  \item an ember, smoke, a flame, to set alight, to be scalded
\end{itemize}

\begin{tcolorbox}[breakable,colback=teal!4,colframe=teal!35!black,title={Before you move on}]
An accident? (\textbf{\gu{અકસ્માત}}, \emph{akasmāt}.) A collision? (\textbf{\gu{ટક્કર}}, \emph{ṭakkar}.) Injured? (\textbf{\gu{ઘાયલ}}, \emph{ghāyal}.) Serious? (\textbf{\gu{ગંભીર}}, \emph{gambhīr}.) To panic? (\textbf{\gu{ગભરાવું}}, \emph{gabhrāvũ}.) To slip? (\textbf{\gu{લપસવું}}, \emph{lapasvũ}.) A graze? (\textbf{\gu{ઉઝરડો}}, \emph{ujharḍo}.) To tumble? (\textbf{\gu{ગબડવું}}, \emph{gabaḍvũ}.) A slope? (\textbf{\gu{ઢાળ}}, \emph{ḍhāḷ}.) Consciousness? (\textbf{\gu{ભાન}}, \emph{bhān}.) A bend? (\textbf{\gu{વળાંક}}, \emph{vaḷā̃k}.) A brake? (\textbf{\gu{બ્રેક}}, \emph{brek}.) Careless? (\textbf{\gu{બેદરકાર}}, \emph{bedarkār}.) A fine? (\textbf{\gu{દંડ}}, \emph{danḍ}.) To crush? (\textbf{\gu{કચડવું}}, \emph{kacaḍvũ}.) An ambulance? (\textbf{\gu{એમ્બ્યુલન્સ}}, \emph{embyulans}.) To pick up? (\textbf{\gu{ઉપાડવું}}, \emph{upāḍvũ}.) A bandage? (\textbf{\gu{પાટો}}, \emph{pāṭo}.) A stitch? (\textbf{\gu{ટાંકો}}, \emph{ṭā̃ko}.) A condition? (\textbf{\gu{હાલત}}, \emph{hālat}.) To warn? (\textbf{\gu{ચેતવવું}}, \emph{cetavvũ}.) Careful? (\textbf{\gu{સાવધાન}}, \emph{sāvdhān}.) Upside down? (\textbf{\gu{ઊંધું}}, \emph{ū̃dhũ}.) A gathered crowd? (\textbf{\gu{ટોળું}}, \emph{ṭoḷũ}.) A push? (\textbf{\gu{ધક્કો}}, \emph{dhakko}.) Death? (\textbf{\gu{મૃત્યુ}}, \emph{mṛtyu}.) Mourning? (\textbf{\gu{શોક}}, \emph{śok}.) A shock? (\textbf{\gu{આઘાત}}, \emph{āghāt}.) Consolation? (\textbf{\gu{આશ્વાસન}}, \emph{āśvāsan}.) To bear? (\textbf{\gu{સહેવું}}, \emph{sahevũ}.) A scream? (\textbf{\gu{ચીસ}}, \emph{cīs}.) A shout? (\textbf{\gu{બૂમ}}, \emph{būm}.) To bring to a stop? (\textbf{\gu{રોકવું}}, \emph{rokvũ}.) To hold back? (\textbf{\gu{અટકાવવું}}, \emph{aṭkāvvũ}.) Safe? (\textbf{\gu{સલામત}}, \emph{salāmat}.) An ember? (\textbf{\gu{અંગારો}}, \emph{angāro}.) Smoke? (\textbf{\gu{ધુમાડો}}, \emph{dhumāḍo}.) A flame? (\textbf{\gu{જ્વાળા}}, \emph{jvāḷā}.) To set alight? (\textbf{\gu{સળગાવવું}}, \emph{saḷgāvvũ}.) To be scalded? (\textbf{\gu{દાઝવું}}, \emph{dājhvũ}.)
\end{tcolorbox}
